% !TeX root = ../paper.tex

\section{讨论}

\subsection{伙伴上下文提供了信息，但其作用不是局部一致改进}

FlexPPI 在 ITrace 全结构上的平均排序优于逐链 FastProtFlex，也优于同一检查点的
孤立链输入。这一结果说明，多链结构中确实存在单链模型无法利用的信息。分层分析
同时表明，该信息主要改善非界面和整体结构范围的排序，而不是直接提高接触残基的
预测准确性。完整输入在非界面的 Spearman 增益为正，在界面则为负；绝对 MAE 也
随伙伴加入略有上升。因此，“模型利用了复合物上下文”和“模型正确学习了界面
动力学”是两个不同命题，整体相关系数只能支持前者。

固定坐标消融解释了这种分离。加入伙伴后，模型对几乎所有非肽界面残基都施加负向
RMSF 修正，并在肽中央 P4--P8 同步降低 RMSF、$H_\chi$ 和 SASA。该响应与重原子
接触距离高度相关，说明局部几何已进入表示；但模型将“接触导致受限”近似为单向
规则，没有充分表示界面两侧的链别差异。对 MHC 接触面，这一规则与 MD 中的相对
刚性一致；对 TCR$\beta$ 接触面及 CDR1/2$\beta$，相同规则却抹除了真实的相对
柔性。该结果提示，下一代模型不应只加强跨链注意力，还需要能够在相同接触几何下
区分伙伴角色、残基化学环境和潜在增柔机制。

\subsection{残基级外部评估比群体平均模式更严格}

九肽群体均值的 Pearson 达到 0.684，容易形成模型已恢复肽动力学图景的印象；然而，
逐复合体平均 Spearman 仅为 0.302，最活跃残基命中率为 34.0\%。二者的差异来自
评价粒度：群体均值主要检验模型是否学到肽位置的平均先验，逐复合体指标则要求模型
识别同一位置在不同 TCR--pMHC 环境中的变化。类似地，模型预测的接触残基在队列
平均上更刚，但这一方向在单个复合体中只略高于随机比例。对需要选择突变位点或实验
测量位点的应用，逐体系残基定位比群体均值更接近实际决策目标，应被视为主要证据。

7L1D/7RRG 个案和氨基酸类型残差可用于生成后续假设，但不能提高外部准确性的证据
等级。前者证明同一肽中央段响应可在两个结构背景中内部复现，后者显示侧链构象熵
比主链 RMSF 更依赖残基类型；两项分析均缺少配对 $H_\chi$/SASA 真值。本文因而
将它们保留为模型可解释性结果，而不表述为已经验证的 TCR 识别机制。

\subsection{上下文响应、结合响应与结合自由能的边界}

从完整复合物中删除伙伴会同时改变跨链边、输入长度和全局注意力的归一化集合；保留
结构仍处于结合态构象。由此得到的差值反映模型对输入上下文的敏感性，而不是蛋白质
从游离态到结合态的物理路径。抗体 CDR 的结构分析已表明，局部环构象改变、环相对
框架移动和构象选择需要分别定义\cite{liu2024cdr}；同理，真实动态响应需要两态
独立结构及轨迹，不能由删链操作替代。

逐残基柔性变化也不能直接换算为结合亲和力。标准结合自由能取决于结合态与游离态
自由能之差，并共同包含相互作用焓、溶剂重排、构象集合、平移和转动熵等贡献。
RMSF 不包含残基间运动协方差，$H_\chi$ 只覆盖离散化的侧链扭转角，SASA 也不区分
被埋藏表面的化学组成。现有热力学约束模型必须显式比较结合/游离状态或完整热力学
循环\cite{jiao2024boltzmann}。因此，FlexPPI 的三类输出可以作为亲和力模型的候选
特征，但不能通过“刚化越强、亲和力越高”的固定规则得到 $\Delta G$。

\subsection{局限性}

本研究存在五项主要限制。第一，ITrace 主队列排除了精确 PDB 重叠，但尚未对全部
训练链完成 MMseqs2 序列近同源和 Foldseek/CATH/ECOD 结构近同源审计；当前证据
不能上升为新折叠泛化。第二，235/245 个 ITrace 复合体超过预训练长度上限 768，
完整输入和孤立输入又处于不同长度范围，因而上下文差值可能包含长度外推效应。
第三，ITrace 仅为 RMSF 提供外部真值，模型的 $H_\chi$ 与 SASA 响应尚未得到同一
体系的 MD 验证。第四，配对响应数据仅有 37 对，全部进入微调，7 个验证体系还用于
检查点选择；该实验没有独立测试集，也未覆盖多个随机种子。第五，当前比较集中于
FastProtFlex，尚缺少 GNM、ANM、CABS-flex、Flexpert-Seq/Flexpert-3D 等不同原理
方法在完全一致输入和残基集合上的系统基准。

\subsection{后续工作}

后续验证应优先解决证据独立性，而非继续增加模型复杂度。第一步是在全部训练数据与
ITrace 之间完成序列和结构去同源，并分别报告“新序列但旧折叠”与“新折叠”结果；
同时进行长度分层并补充物理和学习型基线。第二步是建立从未参与训练或轮次选择的
apo/holo 多重复轨迹测试集，同时计算 $\Delta$RMSF、$\Delta H_\chi$ 和
$\Delta$SASA；响应头、状态相减和零变化基线应在多随机种子下比较。第三步是在响应
微调中加入绝对状态 replay、多目标检查点选择和角色感知的跨链表示，以减少 SASA
遗忘和单向刚化偏置。

只有在上述状态与响应能力被独立验证后，才适合训练复合物级亲和力头。该阶段应以
真实 $K_d/K_i$ 转换的 $\Delta G_{\mathrm{bind}}^{\circ}$ 或成对
$\Delta\Delta G$ 为监督，按复合物家族和界面模板隔离数据，并直接比较“仅静态
界面特征”与“静态特征加 FlexPPI 动态表示”。只有后者在隔离测试集和多个随机种子
上稳定改进，才能将柔性表示解释为亲和力预测的增量信息。
